\chapter{数词: \textcolor{yellow}{die} Zahl}
\section{基数}
\subsection{基本词}
\tikzset{
    label/.style={
            anchor=center,
        },
}
\begin{center}
    \begin{tikzpicture}
        \node[tight, draw=none, align=center] (t1){
            \begin{tabular}{|c|c|}
                \rowcolor{green}
                \hline
                \# & numeral \\
                \hline
                0  & null    \\
                1  & ein     \\
                2  & zwei    \\
                3  & drei    \\
                4  & vier    \\
                5  & fünf    \\
                6  & sechs   \\
                7  & sieben  \\
                8  & acht    \\
                9  & neun    \\
                10 & zehn    \\
                11 & elf     \\
                12 & zwölf   \\
                \hline
            \end{tabular}
        };

        \node[tight, draw=none, align=center, right=1em of t1.north east, anchor=north west] (t2){
            \begin{tabular}{|c|c|}
                \rowcolor{gray}
                \hline
                \# & numeral  \\
                \hline
                16 & sechzehn \\
                17 & siebzehn \\
                18 & achtzehn \\
                \hline
            \end{tabular}
        };

        \node[tight, draw=none, align=center, below=1em of t2] (t3){
            \begin{tabular}{|c|c|}
                \rowcolor{yellow}
                \hline
                \# & numeral \\
                \hline
                20 & zwanzig \\
                30 & dreißig \\
                60 & sechzig \\
                70 & siebzig \\
                80 & achtzig \\
                \hline
            \end{tabular}
        };

        \node[tight, draw=none, align=center, below=1em of t3] (t4){
            \begin{tabular}{|c|c|}
                \rowcolor{green}
                \hline
                \#  & numeral             \\
                \hline
                100 & hund\underline{ert} \\
                1K  & tausend             \\
                \hline
            \end{tabular}
        };

        \node[tight, draw=none, align=center, right=1em of t2.north east, anchor=north west] (t5){
            \begin{tabular}{|c|c|}
                \rowcolor{green}
                \hline
                \#  & numeral    \\
                \hline
                1M  & die Million    \\
                1B  & die Milliarde \\
                1KB & die Billion    \\
                1MB & die Billiarde  \\
                \hline
            \end{tabular}
        };
    \end{tikzpicture}
\end{center}

\subsection{规则}
德语数词通过\textcolor{green}{简单}，\textcolor{yellow}{派生}，\textcolor{gray}{复合}产生。
\begin{center}
    \begin{tabular}{|l|l|l|}
        \rowcolor{gray}
        \hline
        range       & rule                     & exception          \\
        \hline
        1\char`\~12 & 固定                       &                    \\
        \hline
        十几          & 个+zehn                   & 16, 17, 18         \\
        \hline
        整十          & 十+zig                    & 20, 30, 60, 70 ,80 \\
        \hline
        几十几         & 个+und+十+zig              &                    \\
        \hline
        百至百万        & \#hundert/tausend+und+\# &                    \\
        \hline
        百万以上        & 不再复合, 单位当作名词             &                    \\
        \hline
        ...1        & eins\textrightarrow{ein} &                    \\
        \hline
        一百/千开头      & ein可省                    &                    \\
        \hline
    \end{tabular}
\end{center}

注：
\begin{itemize}
    \item sechs 是k的音，sechzehn,sechzig 是ich-laut
    \item 注意hundert发音\freesans{hʊndɐt}
    \item 小于百万，都是复合词（不分开，无复数
    \item 大于百万，当作名词看（首大写，加量词）
    \item 只有十几用und连接
    \item 百千的6,7,8无特殊
\end{itemize}

\subsection{大数、小数}
\begin{itemize}
    \item 四位以上加千分号 der Punkt: . \\
          e.g.: 123.456.789 = hundertdreiundzwanzig \textcolor{red}{Millionen} \\
          viewhundertsechsundfünfzigtausendsiebenhundertneunundachtzig
    \item 小数点 das Komma: , \\
          e.g.: 1,5 = eins Komma fünf （注意这里是eins，只有在xxx1时候eins变成ein）
\end{itemize}
\begin{myquote}
    大数千分号，小数小数点，英德正相反。
\end{myquote}

\section{量词}
\label{sec:num:unit}
辨析：数量单位，就是物理量单位，计量单位就是量词。 \\
用法：量词+名词，单复看量词。阳中量词无复数(Glas 特殊)，只有阴性量词变复数。名词若可数要变复数。e.g.: 100 Gramm Käse kosten 99 Cent. \\
德语量词较少，多用不定冠。 \\
\subsection{数量单位 \textcolor{yellow}{die} Mengeneinheit}
\begin{center}
    \begin{tabular} {|c|c|c|c|c|}
        \hline
        \rowcolor{gray}
        word       & gender                  & genitive & plural & meaning    \\
        \hline
        Gramm      & \textcolor{purple}{das} &          &        & gram       \\
        Kilo       & \textcolor{purple}{das} &          &        & kilogram   \\
        Kilogramm  & \textcolor{purple}{das} &          &        & kilogram   \\
        Pfund      & \textcolor{purple}{das} &          &        & pound      \\
        \hline
        Liter      & \textcolor{green}{der}  &          &        & liter      \\
        Milliliter & \textcolor{green}{der}  &          &        & milliliter \\
        \hline
    \end{tabular}
\end{center}
注：
\begin{itemize}
    \item 数量单位不用复数。
    \item 换算：Ein Kilo hat zwei Pfund.
\end{itemize}

\subsection{计量单位 \textcolor{yellow}{die} Mengenangabe}
\begin{center}
    \begin{tabular} {|c|c|c|c|c|}
        \hline
        \rowcolor{gray}
        word     & gender                  & genitive & plural & meaning \\
        \hline
        Tasse    & \textcolor{yellow}{die} &          &        & cup     \\
        Flasche  & \textcolor{yellow}{die} &          &        & bottle  \\
        Dose     & \textcolor{yellow}{die} &          &        & can     \\
        Kanne    & \textcolor{yellow}{die} &          &        & pot     \\
        Packung  & \textcolor{yellow}{die} &          &        & pack    \\
        Tafel    & \textcolor{yellow}{die} &          &        & bar     \\
        \hline
        Stück    & \textcolor{purple}{das} &          &        & block   \\
        Glas     & \textcolor{purple}{das} &          &        & glass   \\
        Kännchen & \textcolor{purple}{das} &          &        & jug/pot \\
        \hline
        Kasten   & \textcolor{green}{der}  &          &        & case    \\
        Beutel   & \textcolor{green}{der}  &          &        & bag     \\
        Becher   & \textcolor{green}{der}  &          &        & cup     \\
        \hline
    \end{tabular}
\end{center}

\label{sec:num:unit}
注：
\begin{itemize}
    \item 冠词限定量词（而不是量词后的名词）。若是计量单位，变位根据单位，若是数量单位，动词通常按复数变位。
    \item 物质名词不可数必须要量词，一些特例可以省略量词：zwei Kohl, zwei Bier, zwei Brot...
    \item 可数名词可不用量词
\end{itemize}
为何量词有性？
\begin{itemize}
    \item 可独立做名词
    \item 阳中量词无复数，只有阴性量词变复数
    \item 量词+名词搭配，冠词限定量词
\end{itemize}

\section{电话号码}
两位数两位数读

\section{货币}
读法规则:
\begin{itemize}
    \item 货币单位不变复数
    \item 先整后小，可全可省。如：4,30 Euro
          \begin{itemize}
              \item vier Euro (und) dreißig Cent
              \item vier Euro dreißig (可看作把小数点读作Euro)
          \end{itemize}
    \item 零整读分。如：0,70 Euro = siebzig Cent
\end{itemize}

\section{计算： \textcolor{purple}{das} Rechnen}
\begin{center}
    \begin{tabular}{|c|l|}
        \rowcolor{gray}
        \hline
        operation  & spelling     \\
        \hline
        +          & plus         \\
        -          & minus        \\
        \texttimes & mal          \\
        \textdiv   & durch        \\
        =          & ist (gliech) \\
        \hline
    \end{tabular}
\end{center}

\section{时间 \textcolor{yellow}{die} Uhrzeit}
\subsection{问时间}
\begin{itemize}
    \item Q: Wie spät ist es bitte?
    \item Q: Wie viel Uhr ist es bitte?
    \item A: Es ist xxx.
\end{itemize}
注意Uhr用wie viel (how much)\\

\subsection{问何时}
\begin{itemize}
    \item Q: Wann gehst du zum Supermarkt?
    \item Q: Um wie viel Uhr gehst du zum Supermarkt?
    \item A: Um xxx.
\end{itemize}

\subsection{时间写法}
用点号Punkt分隔，e.g. 15.15
\subsection{时间读法}
null Uhr = 24 Uhr
\subsubsection{offiziell/formell: 24 std.}
A:B = A Uhr B\\
if A = 1, eins \textrightarrow{ein} (e.g. 1:01 = ein Uhr eins) \\
\subsubsection{inoffiziell/informell: 12 Std.}
\begin{center}
    \clock[hand=false]
\end{center}
12小时制除整点同24，用Uhr，其余不用。 \\
上下左右特殊，有专用读法。 (注意halb后的数字比当前点数大1。8:30 = halb 9)\\
其余时间 = \# vor/nach \#/halb \# \\
\subsubsection{Estimation}
\begin{itemize}
    \item 约: gegen. e.g. Es ist gegen eins.
    \item 略: kurz. e.g. Es ist kurz vor/nach eins.
    \item 整: Punkt. e.g. Es ist Punkt eins.
\end{itemize}
注意der Punkt \sqr[color=green]

\subsection{Misc}
带半：
\begin{itemize}
    \item 半个： \\
          einhalb Stunde (num. + n.) =  \\
          eine halbe Stunde (art. + adj. + n.)
    \item 一个半： \\
          eineinhalb Stunden = \\
          anderthalb Stunden (这里用anderthalb可能因为einein太别扭了) \\
    \item N个半：Neinhalb Stunden
\end{itemize}